\section{XGBoost：从假说挑战到行业标准}

\subsection{诞生背景}

2013 年 AlexNet 横空出世后，陈天奇和导师 Carlos Guestrin 讨论：深度学习成功的原因到底是神经网络本身，还是大 hypothesis space + 大数据？如果后者是关键，那树模型也有 arbitrary grow 的能力（不断增加树的数量和深度），或许也能在大数据上取得突破。

他们决定用树模型来挑战这个假说。陈天奇基于 KDD Cup 的经验，坚持选择 tree-based method（而非导师偏好的 kernel method）。他认为树模型有一个很强的特性：可以随着数据量 arbitrary grow depth，而 kernel method 虽然也是 nonparametric，但在 tabular data 上没有同样的优势。Carlos 给了他充分的自由度------``我非常感谢我的导师，他给了我非常大的自由度。''

XGBoost 的名字是 eXtreme Gradient Boosting，``extreme'' 体现了做到极致的追求。

\begin{importantbox}{XGBoost 成功的三个关键}
\begin{enumerate}
    \item \textbf{极致性能}：在发布时是世界上最快的 gradient boosting 实现。陈天奇把每一个优化都做到极致，这源自 ACM 班``做一件事情要做到极致''的传统。
    \item \textbf{开源社区}：XGBoost 的成功很大程度来自社区贡献。例如一位法国律师贡献了 R 语言可视化包，用于检测税务欺诈。社区力量让项目远超个人能力范围。
    \item \textbf{专注}：只做一个算法，把这一个算法做好做到极致。相比之前的``大而全''项目 SVD Feature，XGBoost 的专注策略更适合小团队。
\end{enumerate}
\end{importantbox}

\subsection{算法与系统的垂直整合}

XGBoost 内置了对 missing value 的自动处理：如果某个特征值缺失，模型会根据历史 pattern 自动决定如何分类，不需要用户预先做数据清洗。这个功能``非常非常受欢迎''，因为实际数据中缺失值无处不在。传统做法是用中位数等方式填充，但 XGBoost 让模型自己学习最优的处理方式。

陈天奇在做 XGBoost 时第一次意识到：\textbf{做机器学习系统，只做算法或只做系统都不够，必须垂直整合}。这个认知贯穿了他此后所有项目。他指出，XGBoost 的成功不仅因为系统层面跑得快（多线程、out-of-core learning、分布式计算），也因为算法层面的细节优化（正则化、sparsity-aware split finding、missing value handling）。\textbf{只有同时理解算法的需求和系统的约束，才能做出真正好用的工具}。

\begin{knowledgebox}{XGBoost 论文的写作哲学}
XGBoost 的论文直到 PhD 第三年才写。Carlos 的要求不是追求中稿，而是``写一篇让读者可以学到东西的文章''。同年 KDD 上 Carlos 实验室还有另一篇 LIME（Explainable ML），由 Marco Ribeiro 完成，开创了可解释机器学习这个方向。Carlos 的研究 style 是鼓励学生做能开创一个领域的工作，而不是在已有方向上做增量改进。
\end{knowledgebox}

\subsection{树模型 vs 神经网络的反思}

十多年后回顾，他认为当时的判断``partly right, partly wrong''：

\begin{itemize}
    \item \textbf{正确的部分}：模型 capacity 不应受限制。树模型可以 arbitrary grow depth，这个特性很强。在 tabular data 和 time series 上，树模型至今仍是 go-to tool。
    \item \textbf{没有预料到的}：神经网络通过一种``更暴力''的方式实现了 capacity 扩展------不是逐渐 grow parameter space，而是一开始就把 parameter space 设得非常大。
    \item \textbf{神经网络的组合特性}：transformer block、linear block 等模块可以相互组合，使得表达能力通过组合更 compact 地表示。
\end{itemize}

\subsection{导师 Carlos Guestrin 的影响}

Carlos 对学生要求极高：第一次见面就说``You should write the best paper in the world''。他的实验室有``no beamer, only PowerPoint''的规定，强调用视觉化方式让读者快速理解研究内容。XGBoost 论文直到博士第三年才写，Carlos 的要求是``写一篇让读者可以学到东西的文章''，而不是优化中稿率。

\subsection{本章小结}

XGBoost 的成功不是偶然：极致的性能、活跃的社区、专注的策略，加上算法与系统的垂直整合，使它成为引用近 7 万次的行业标准。但陈天奇当时并没有预料到这个成功------纯粹是``做了一个我们觉得还行的东西，然后开源了''。

